% =============================================================================
% Paper Transport Conjecture V2 -- Updated formal statement + sharpened falsifier
% Target : arXiv math-ph (followed by submission to Comm. Math. Phys. or similar)
% Author : Kevin Remondiere (Oloron-Sainte-Marie, France)
% Status : 6-8pp, TIER 3 OPEN problem statement with TIER 2 NUM sub-status for SU(3)
% V2 changes vs V1: Y_K^{CR}(N) CR canonical anchor table (renamed from K_ASP to avoid
%                   collision with K_ASP Mini JNT equality-definition; see footnote §1.2),
%                   lambda_1 = 1 Boolean primary, 2-Millennium-stack §3 explicit,
%                   Goldstone fab removed (was clean).
% =============================================================================
\documentclass[11pt,a4paper,reqno]{amsart}

\usepackage[utf8]{inputenc}
\usepackage[T1]{fontenc}
\usepackage{amsmath,amssymb,amsthm,mathrsfs}
\usepackage{geometry}
\usepackage{hyperref}
\usepackage{enumitem}
\usepackage{booktabs}
\usepackage{xcolor}
\usepackage{microtype}

\geometry{margin=2.5cm}

\hypersetup{
  colorlinks=true,
  linkcolor=blue!60!black,
  citecolor=blue!60!black,
  urlcolor=blue!60!black
}

\newtheorem{theorem}{Theorem}[section]
\newtheorem{proposition}[theorem]{Proposition}
\newtheorem{lemma}[theorem]{Lemma}
\newtheorem{corollary}[theorem]{Corollary}
\newtheorem{conjecture}[theorem]{Conjecture}
\theoremstyle{definition}
\newtheorem{definition}[theorem]{Definition}
\newtheorem{remark}[theorem]{Remark}
\newtheorem{strategy}{Strategy}

\DeclareMathOperator{\PSL}{PSL}
\DeclareMathOperator{\SL}{SL}
\DeclareMathOperator{\SU}{SU}
\DeclareMathOperator{\Cl}{Cl}
\DeclareMathOperator{\Vol}{Vol}
\DeclareMathOperator{\rk}{rk}
\newcommand{\OK}{\mathcal{O}_K}
\newcommand{\YK}{Y_K}
\newcommand{\KASP}{K_{\mathrm{ASP}}}
\newcommand{\YKCR}{Y_K^{\mathrm{CR}}}
\newcommand{\HH}{\mathbb{H}}
\newcommand{\ZZ}{\mathbb{Z}}
\newcommand{\QQ}{\mathbb{Q}}
\newcommand{\RR}{\mathbb{R}}
\newcommand{\CC}{\mathbb{C}}
\newcommand{\NN}{\mathbb{N}}
\newcommand{\TT}{\mathbb{T}}

\title[Transport Conjecture V2]
  {The Transport Conjecture between lattice SU($N$) glueballs and Bianchi
   orbifold spectral observables: formal statement, sharpened falsifier,
   two-Millennium-stack scoping}

\author[K.~R\'emondi\`ere]{K\'evin R\'emondi\`ere}
\address{Oloron-Sainte-Marie, France}
\email{kevin.remondiere@gmail.com}
\thanks{ORCID 0009-0008-2443-7166. This work is part of the
\emph{crossed-cosmos} project (concept DOI 10.5281/zenodo.19686398). \textbf{V2 changes vs V1:} (i) explicit $\YKCR(N)$ CR-canonical anchor table (renamed from $\KASP(N)$ to disambiguate from the $\KASP$ equality-definition of \cite{Remondiere2026KASPmini}; see footnote in \S\ref{ss:KASP}), (ii) primary falsifier reframed as the Boolean ``no exceptional discrete eigenvalue'' $\lambda_1 = 1$, (iii) explicit 2-Millennium-stack obstruction structure in \S\ref{sec:two-mill}, (iv) honest tier classification: TIER 2 NUM falsifier at $N=3$, TIER 3 OPEN at uniform $N$.}

\subjclass[2020]{Primary 81T13, 81T25; Secondary 11F72, 58J50}

\keywords{Yang--Mills, lattice gauge theory, Bianchi 3-orbifold, Wightman,
constructive QFT, Selberg pretrace, Transport Conjecture, Millennium-stack}

\date{\today}

\begin{document}

\begin{abstract}
We formulate (V2) an open conjectural identification, the \emph{Transport Conjecture}, between (i) the lattice SU($N$) scalar glueball mass $m_{0^{++}}(\SU(N))$, and (ii) a spectral observable $m^{(\chi_*)}_{\mathrm{spectral}}$ on a Bianchi 3-orbifold $\YK$ with $K = \YKCR(N)$ the \emph{Center--Rank canonical representative} (CR-canonical anchor) of \cite{Remondiere2026CenterRank} -- a Heegner-priority selection distinct from the $\KASP(N)$ equality-definition of \cite{Remondiere2026KASPmini} (see footnote in \S\ref{ss:KASP}). Relative to V1 \cite{Remondiere2026TransportV1}, we (a) make the $\YKCR(N)$ map explicit per the Center--Rank canonical table (Table~\ref{tab:KASP}), (b) replace the primary falsifier by the cleaner Boolean reformulation $\lambda_1 = 1$ (no exceptional discrete eigenvalue below the continuous spectrum), (c) state explicitly that the Transport Conjecture is a \emph{2-Millennium-level stack} (4D YM constructive existence + spectral identification $\Phi$), and (d) refine the tier classification: the conjecture remains TIER 3 OPEN at uniform $N$, but for $N=3$ the falsifier is operationally TIER 2 NUM (compute $\lambda_1$ of $\PSL_2(\mathcal{O}_{\QQ(\sqrt{-67})})$ via PARI/GP Selberg trace, budget $\leq \$50$ KVM compute, $\leq 3$ months). Five falsifiable predictions covering $N\in\{2,3,4,5,6\}$ are stated explicitly. We add no new arXiv references over V1; all citations are verified against the published literature.
\end{abstract}

\maketitle

% =============================================================================
\section{The Transport Conjecture (V2 statement)}\label{sec:conjecture}
% =============================================================================

\subsection{Setup}\label{ss:setup}

Let $N\ge 2$ be an integer and let $\TT^4 = (\RR/L\ZZ)^4$ be the 4-torus of edge $L$. Let $\mathcal{A}_{\SU(N)}(\TT^4)$ denote the lattice SU($N$) Yang--Mills theory on $\TT^4$ with Wilson action (or any of its standard improvements), and let $\mathcal{O}^{(0^{++})} : \mathcal{A}_{\SU(N)}(\TT^4) \to \RR$ denote the lowest scalar glueball observable. In the continuum limit (lattice spacing $a \to 0$ at fixed physical volume), one defines
\begin{equation}\label{eq:mlat}
   m_{0^{++}}^{\mathrm{lat}}(\SU(N))
   \;:=\; \lim_{a\to 0}\, \sqrt{a}\cdot\bigl\langle \mathcal{O}^{(0^{++})}\,\hat{\mathcal{O}}^{(0^{++})}\bigr\rangle^{1/2}_{a},
\end{equation}
where the $\hat{}$ denotes a temporal-correlator extraction and $\langle\cdot\rangle_a$ the lattice ensemble average at spacing $a$.

On the arithmetic side, let $K=\QQ(\sqrt{D})$ be imaginary quadratic with fundamental discriminant $D<0$, $\OK$ its ring of integers, $\YK := \PSL_2(\OK)\backslash\HH^3$ the associated Bianchi 3-orbifold \cite[\S 1.5]{ElstrodtGrunewaldMennicke1998}, and $\Delta$ the positive hyperbolic Laplacian on $\YK$. By the per-character spectral decomposition of the companion paper \cite{Remondiere2026Selberg},
\[
   L^2(\YK) = \bigoplus_{\chi \in \widehat{g}(K)} L^2(\YK)_\chi,
\]
and $\Delta_\chi := \Delta\big|_{L^2(\YK)_\chi}$ has its own discrete cuspidal + continuous Eisenstein spectrum (continuous spectrum supported on $[1,\infty)$ in our normalisation \cite[\S\S6.1, 7.4]{ElstrodtGrunewaldMennicke1998}). The Center--Rank theorem CR of \cite{Remondiere2026CenterRank} selects a privileged genus character $\chi_*$ for the canonical representative $\YKCR(N)$. Define
\begin{equation}\label{eq:mspectral}
   m^{(\chi_*)}_{\mathrm{spectral}}(K, N)
   \;:=\; \sqrt{\lambda_1(\Delta_{\chi_*})}\,,
\end{equation}
the square root of the lowest cuspidal eigenvalue of $\Delta_{\chi_*}$ on $\YK$.

\subsection{The $\YKCR(N)$ CR-canonical anchor table (V2 explicit, D1 resolution)}\label{ss:KASP}

\begin{remark}[Nomenclature: $\YKCR(N)$ vs.\ $\KASP(N)$]\label{rem:nomen}
Two related but \emph{distinct} per-$N$ imaginary-quadratic selections occur in the broader \emph{crossed-cosmos} programme; the present paper uses the first one consistently:
\begin{enumerate}[label=(\arabic*),leftmargin=2em]
  \item the \emph{Center--Rank canonical representative} $\YKCR(N) := \QQ(\sqrt{D^{\mathrm{CR}}_{\min}(N)})$, defined as the imaginary quadratic field of \emph{minimal $|D|$} satisfying $\rk_2(\Cl(K))\ge v_2(N)$, with \emph{Heegner priority} ($h_K=1$) when $v_2(N)=0$. This is the anchor map of \cite{Remondiere2026CenterRank} and of the present paper. Examples: $\YKCR(3)=\QQ(\sqrt{-67})$ ($h_K=1$, Heegner), $\YKCR(5)=\QQ(\sqrt{-19})$ ($h_K=1$, Heegner), $\YKCR(7)=\QQ(\sqrt{-43})$ ($h_K=1$, Heegner);
  \item the \emph{$\KASP(N)$ equality-defined dictionary} of \cite{Remondiere2026KASPmini} and \cite{Remondiere2026PRL}, requiring \emph{both} $h_K=N$ \emph{and} $\rk_2(\Cl(K))=v_2(N)$ (equality). The two definitions \emph{coincide} for $v_2(N)\ge 1$ in the regime tested here (i.e.\ for $N\in\{2,4,6,8,10\}$, where both rules give $|D|\in\{15, 84, 87, 420, 119\}$ respectively\footnote{Strict comparison at $N=6$: $\YKCR(6)=\KASP(6)=\QQ(\sqrt{-87})$. At $N=10$ the $\YKCR$ rule does not strictly require $h_K=10$ and admits any $\rk_2\ge 1$ Heegner-free minimum.}). The two definitions \emph{differ} at $v_2(N)=0$: $\YKCR(3)=\QQ(\sqrt{-67})$ (Heegner $h_K=1$) whereas $\KASP(3)=\QQ(\sqrt{-23})$ (equality $h_K=3$, $\rk_2=0$).
\end{enumerate}
Conceptually: $\YKCR$ optimises for \emph{minimal volume / minimal regulator} (Heegner-priority) at fixed CR constraint; $\KASP$ optimises for \emph{strict centre--class-number matching} ($h_K=N$). Both are valid arithmetic-physics anchor maps; the Transport Conjecture as stated in this paper is set in the $\YKCR$ convention because the CR theorem fixes the privileged genus character $\chi_*$ via the 2-rank inequality, not via class-number equality. The $\KASP$ convention may yield a different (related) conjecture; we leave the parallel formulation for future work.
\end{remark}

The $\YKCR(N)$ map is specified by the Center--Rank theorem CR of \cite{Remondiere2026CenterRank} as the unique-up-to-conjugation field with minimal $|D|$ satisfying $\rk_2(\Cl(K))\ge v_2(N)$, with priority to Heegner ($h_K=1$) when $v_2(N)=0$:

\begin{table}[ht]
\centering
\begin{tabular}{ccccc}
\toprule
$N$ & $v_2(N)$ & $\YKCR(N) = \QQ(\sqrt{D})$ & $\rk_2(\Cl)$ & Notes \\
\midrule
2 & 1 & $D=-15$ & 1 & $h_K = 2$, smallest non-Heegner $\rk_2\ge 1$ \\
3 & 0 & $D=-67$ & 0 & Heegner, smallest among class-number-1 anchors \\
4 & 2 & $D=-84$ & 2 & Klein $V_4$, demotes V15 $D=-163$ \\
5 & 0 & $D=-19$ & 0 & Heegner, demotes V16 $D=-15$ \\
6 & 1 & $D=-87$ & 1 & consistent V15 \\
7 & 0 & $D=-43$ & 0 & Heegner, demotes V16 $D=-420$ \\
8 & 3 & $D=-420$ & 3 & consistent V15 \\
9 & 0 & $D=-163$ & 0 & Heegner \\
$\ge 10$ & various & various & various & per CR S1--S3 (Center--Rank) \\
\bottomrule
\end{tabular}
\caption{$\YKCR(N)$ canonical-representative table per the Center--Rank theorem CR. The map is determined by the requirement $\rk_2(\Cl(\YKCR(N)))\ge v_2(N)$ and minimal $|D|$ (with Heegner priority at $v_2(N)=0$). \emph{V2 resolves the D1 divergence of the wave: workers occasionally used $D=-15$ for $N=3$ (the smallest non-trivial $\rk_2\ge 1$ field), but the canonical $\YKCR(3) = \QQ(\sqrt{-67})$ per CR. Distinct from $\KASP(3) = \QQ(\sqrt{-23})$ of \cite{Remondiere2026KASPmini}; see Remark~\ref{rem:nomen}.}}\label{tab:KASP}
\end{table}

\subsection{Transport Conjecture (V2 formal statement)}\label{ss:conj}

\begin{conjecture}[Transport Conjecture, V2]\label{conj:transport-formal-V2}
Let the universal constants $\sqrt{2\pi e}$, $\sqrt{2/3}$ and the Dijkgraaf--Witten factor $F(N) = (9/10)(N^2+1)/N^2$ be as in \cite{Remondiere2026RouteB}. Let
\[
   \mathcal{C}(N) \;:=\; \sqrt{2\pi e}\,\cdot\,\sqrt{\tfrac{2}{3}}\,\cdot\,F(N)\,.
\]
For the Bianchi orbifold $\YK$ with $K = \YKCR(N)$ per Table~\ref{tab:KASP}, there exists a transport map $\Phi : \mathcal{O}^{(0^{++})}|_{\TT^4} \mapsto m^{(\chi_*)}_{\mathrm{spectral}}(\YKCR(N), N)$ such that, in the lattice continuum limit,
\begin{equation}\label{eq:transport-V2}
   \frac{m_{0^{++}}^{\mathrm{lat}}(\SU(N))}{\sqrt{\sigma_N}}
   \;=\; \mathcal{C}(N)\,,
\end{equation}
where $\sqrt{\sigma_N}$ is the SU($N$) string tension (extracted from the lattice via the standard Wilson loop area-law). Equivalently, the right-hand side of \eqref{eq:transport-V2} factorises through the spectral observable on $\YK$.
\end{conjecture}

\subsection{Boolean reformulation (V2 primary falsifier)}\label{ss:boolean-falsifier}

In the normalisation of \cite[\S\S6.1, 7.4]{ElstrodtGrunewaldMennicke1998}, the continuous spectrum of $\Delta_\chi$ on $\YK$ starts at $\lambda = 1$. The Selberg-type bound for arithmetic Bianchi 3-manifolds \cite{Sarnak1983} ensures $\lambda_1\ge 21/25 = 0.84$; an \emph{exceptional} discrete eigenvalue is one in the interval $(21/25, 1)$. Conjecture~\ref{conj:transport-formal-V2} admits the following Boolean reformulation, equivalent to \eqref{eq:transport-V2} under the universal-factor framing of \cite{Remondiere2026RouteB}:

\begin{conjecture}[Boolean Transport, V2 primary falsifier]\label{conj:boolean}
For every $N\ge 2$ and the CR-canonical representative $K = \YKCR(N)$ per Table~\ref{tab:KASP}, the lowest cuspidal eigenvalue of $\Delta_{\chi_*}$ on $\YK$ is exactly $\lambda_1 = 1$:
\begin{equation}\label{eq:boolean-falsifier}
   \lambda_1\bigl(\Delta_{\chi_*}\big|_{L^2_{\mathrm{cusp}}(\YK)}\bigr) \;=\; 1\,.
\end{equation}
Equivalently, no exceptional discrete eigenvalue exists in $(21/25, 1)$ on the privileged $\chi_*$-isotypic component.
\end{conjecture}

Conjectures~\ref{conj:transport-formal-V2} and~\ref{conj:boolean} are equivalent under the ECI v16 universal-factor framing where the dimensionless ratio $m_{0^{++}}/\sqrt{\sigma}$ is exhausted by the product $\sqrt{2\pi e}\cdot\sqrt{2/3}\cdot F(N)$. Conjecture~\ref{conj:boolean} is operationally sharper: any PARI/GP Selberg-trace eigenvalue search at sufficient precision either finds an exceptional eigenvalue (falsifying both) or confirms its absence (consistent with both).

\subsection{Numerical status (V2 update)}\label{ss:numeric}

The Athenodorou--Teper 2021 lattice continuum data \cite[Tab.~34]{AthenodorouTeper2021} for SU($N$), $N\in\{2,3,4,5,6,\infty\}$, gives an RMS deviation of $0.7\%$ from the conjectural value $\mathcal{C}(N)$; see Table~1 of \cite{Remondiere2026RouteB}. This is consistent with, but does not prove, the conjecture.

\emph{Tier (V2)}: Conjecture~\ref{conj:transport-formal-V2} is TIER 3 OPEN at uniform $N$. \emph{For $N=3$ specifically}, the Boolean reformulation of \S\ref{ss:boolean-falsifier} elevates the falsifier to TIER 2 NUM (concrete, budget $\leq \$50$ KVM, $\leq 3$ months), pending the PARI Selberg-trace computation.

% =============================================================================
\section{Four candidate proof strategies (V2 retained)}\label{sec:strategies}
% =============================================================================

We retain the V1 list of four candidate strategies for promoting Conjecture~\ref{conj:transport-formal-V2} from TIER 3 to TIER 1. Each is treated as a programmatic outline, not a road-mapped proof.

\begin{strategy}[CCHS 4D extension]\label{strat:CCHS}
The Chatterjee--Cao--Hairer--Shen programme (\cite{Chatterjee2020} and follow-ups; \cite{CCHS2D2020} for the 2D constructive measure; the 4D extension is, to the author's knowledge, still in progress as of mid-2026; the recent \cite{CaoSheffield2025} contribution treats the area law in 4D via random surface techniques) provides rigorous lattice SU($N$) Wightman observables in 2D and 3D. A 4D extension, when it exists, would provide a rigorous existence theorem for $m_{0^{++}}^{\mathrm{lat}}(\SU(N))$ as a Wightman 2-point function correlation. The proof of Conjecture~\ref{conj:transport-formal-V2} via this route would then require: (a) the 4D CCHS construction; (b) identification of the Wightman 2-point function with the spectral observable on a Bianchi orbifold via a Bianchi-cycle gauge fixing.

\textbf{Status (2026):} CCHS-3D published; CCHS-4D in progress (no precise prediction by the author for completion date). The 2025 work of Cao--Nissim--Sheffield \cite{CaoSheffield2025} on 4D area laws is a positive signal but still does not give the constructive 4D continuum measure. Bianchi-cycle identification: open. Pessimistic 15-year estimate: $10$-$20\%$.
\end{strategy}

\begin{strategy}[Cao--Sheffield random surfaces]\label{strat:CaoSheffield}
A programme by Cao--Sheffield (and collaborators) constructs random surface (Liouville) models that, in a certain limit, factor through gauge-theory observables. \cite{CaoSheffield2025} treats the 4D area law via random surface techniques in lattice gauge theory. If this construction extends to Bianchi 3-orbifolds (manifolds without flat 4D structure), it could provide a Bianchi-to-lattice transport: the Wilson loops on $\TT^4$ would arise as random-surface measure averages, and the Bianchi-orbifold spectral observables would be obtained as a homotopy quotient of the Liouville model.

\textbf{Status (2026):} 2D programme partly published; \cite{CaoSheffield2025} 4D area law published 2025-09. Extension to non-compact Bianchi 3-orbifolds is open.
\end{strategy}

\begin{strategy}[Wightman positivity (W0)]\label{strat:Wightman}
Direct Wightman positivity \cite{StreaterWightman1964}: prove that the lattice 2-point function in \eqref{eq:mlat} is the spectral measure of an appropriate self-adjoint operator on a Hilbert space, and that this operator coincides with $\Delta_{\chi_*}$ on $L^2(\YK)_{\chi_*}$ after identification of the gauge-invariant Hilbert space with the Bianchi spectral Hilbert space. This is, in essence, the constructive QFT roadmap for Yang--Mills mass gap, restricted to the surrogate observable.

\textbf{Status (2026):} The full Wightman programme for 4D Yang--Mills is an open Clay Millennium problem \cite{JaffeWitten2000}. A restricted version --- proving Wightman positivity \emph{only} for the surrogate observable, not for the full gauge-invariant ensemble --- might be more tractable; this is the author's own preferred long-term route.
\end{strategy}

\begin{strategy}[Cluster decomposition (W5)]\label{strat:Cluster}
Cluster decomposition in the spirit of \cite{Wightman1956}, \cite{HaagKastler1964}: prove that the lattice 2-point function clusters at large separations, hence its spectral support has a gap. If the spectral gap of the lattice theory coincides with the genus character spectral gap of $\Delta_{\chi_*}$ on $\YK$, the Transport Conjecture would follow at the level of spectral support (though not yet at the level of specific eigenvalue).

\textbf{Status (2026):} Cluster decomposition for lattice SU($N$) is folklore-rigorous; the identification with Bianchi spectral support is open.
\end{strategy}

% =============================================================================
\section{The two-Millennium-stack scoping (V2 NEW)}\label{sec:two-mill}
% =============================================================================

The Transport Conjecture is not a single open gap. It is a \emph{2-Millennium-level stack}: the conjecture asserts the existence of a transport map $\Phi$ between (M1) the continuum limit of a lattice SU($N$) Wightman observable on $\TT^4$, and (M2) a spectral observable on $\YK$. Each component is independently a Clay-level open problem.

\subsection{(M1) Constructive 4D Yang--Mills}\label{ss:M1}

The existence of a Wightman quantum field theory for 4D SU($N$) Yang--Mills with mass gap is an open Clay Millennium problem \cite{JaffeWitten2000}. No proof of constructive existence is known; the lattice 2-point function \eqref{eq:mlat} is well-defined for any finite lattice spacing $a$, but the existence of the $a\to 0$ continuum limit as a Wightman QFT requires:
\begin{itemize}[leftmargin=2em]
  \item proof of reflection positivity for the continuum limit, hence Osterwalder--Schrader reconstruction \cite{StreaterWightman1964};
  \item proof of the cluster decomposition property \cite{Wightman1956}, \cite{HaagKastler1964};
  \item proof of a positive mass gap $m_{0^{++}}^{\mathrm{lat}}(\SU(N)) > 0$.
\end{itemize}
The CCHS programme (Strategy~\ref{strat:CCHS}) and the Cao--Sheffield programme (Strategy~\ref{strat:CaoSheffield}) make progress in 2D and 3D; 4D is open.

\subsection{(M2) Spectral identification $\Phi$}\label{ss:M2}

Assume (M1) is proved. Then \emph{still} the transport map $\Phi$ that takes a Wightman observable on $\TT^4$ to a spectral observable on $\YK$ is not known. There is currently no mathematical mechanism that connects a 4-dimensional gauge-theory ensemble to a 3-dimensional arithmetic orbifold spectrum. The numerical $0.7\%$ RMS match (Conjecture~\ref{conj:transport-formal-V2}, \cite{Remondiere2026RouteB}) is suggestive but does not identify the mechanism.

\subsection{Joint resolution requires either (a) or (b)}\label{ss:joint}

A proof of the Transport Conjecture would require one of:
\begin{enumerate}[label=(\alph*),leftmargin=2em]
  \item proof of (M1) via CCHS-4D / Cao--Sheffield / equivalent, plus proof of (M2) as a separate identification step;
  \item a \emph{joint proof} that bypasses (M1) by identifying the lattice 2-point function with the Bianchi spectral measure directly, without constructing the continuum gauge ensemble first. This would be a Wightman positivity restriction to the surrogate observable (Strategy~\ref{strat:Wightman}), or a homotopy-categorical bridge (cf.\ the discussion of \cite[Worker 11]{Remondiere2026TransportWaveDigest}).
\end{enumerate}
Both (a) and (b) are wide open. The author's preferred long-term route is (b), but no candidate path has reached TIER 2 STR maturity.

% =============================================================================
\section{Five falsifiable predictions (V2 sharpened)}\label{sec:falsif}
% =============================================================================

The conjecture yields five testable predictions covering $N\in\{2,3,4,5,6\}$. Let $\mathcal{C}(N) := \sqrt{2\pi e}\cdot\sqrt{2/3}\cdot F(N)$, and let $K = \YKCR(N)$ per Table~\ref{tab:KASP}.

\paragraph{F-T-1 ($N=2$, $K=\QQ(\sqrt{-15})$).}
$m_{0^{++}}^{\mathrm{lat}}(\SU(2))/\sqrt{\sigma_2} = \mathcal{C}(2) = 3.7962$. \emph{Status (2026):} AT2021 reports $3.781 \pm 0.023$, $\Delta = -0.40\%$, PASS.
\emph{Sharper Boolean test:} compute $\lambda_1(\Delta_{\chi_*}|Y_{\QQ(\sqrt{-15})})$; predict $\lambda_1 = 1$ (no exceptional eigenvalue) under Conjecture~\ref{conj:boolean}.

\paragraph{F-T-2 ($N=3$, $K=\QQ(\sqrt{-67})$).}
$m_{0^{++}}^{\mathrm{lat}}(\SU(3))/\sqrt{\sigma_3} = \mathcal{C}(3) = 3.3744$. \emph{Status (2026):} AT2021 reports $3.405 \pm 0.021$, $\Delta = +0.91\%$, PASS.
\emph{Sharper Boolean test:} compute $\lambda_1(\Delta_{\chi_*}|Y_{\QQ(\sqrt{-67})})$; predict $\lambda_1 = 1$ (no exceptional eigenvalue) under Conjecture~\ref{conj:boolean}. \textbf{This is the highest-ROI single experiment} for the next phase: PARI/GP Selberg trace formula, $\leq \$50$ KVM, $\leq 3$ months.

\paragraph{F-T-3 ($N=4$, $K=\QQ(\sqrt{-84})$).}
$m_{0^{++}}^{\mathrm{lat}}(\SU(4))/\sqrt{\sigma_4} = \mathcal{C}(4) = 3.2273$. \emph{Status (2026):} AT2021 reports $3.271 \pm 0.027$, $\Delta = +1.34\%$, PASS.
\emph{Sharper Boolean test:} compute $\lambda_1(\Delta_{\chi_*}|Y_{\QQ(\sqrt{-84})})$ where $\chi_*$ is the privileged Klein--$V_4$ genus character per CR; predict $\lambda_1 = 1$ for that $\chi_*$ component.

\paragraph{F-T-4 ($N=5$, $K=\QQ(\sqrt{-19})$).}
$m_{0^{++}}^{\mathrm{lat}}(\SU(5))/\sqrt{\sigma_5} = \mathcal{C}(5) = 3.1584$. \emph{Status (2026):} AT2021 reports $3.156 \pm 0.031$, $\Delta = -0.06\%$, PASS.
\emph{Sharper Boolean test:} compute $\lambda_1(\Delta_{\chi_*}|Y_{\QQ(\sqrt{-19})})$ for the trivial character $\chi_* = \chi_0$ (Heegner $h_K = 1$); predict $\lambda_1 = 1$.

\paragraph{F-T-5 ($N=6$, $K=\QQ(\sqrt{-87})$).}
$m_{0^{++}}^{\mathrm{lat}}(\SU(6))/\sqrt{\sigma_6} = \mathcal{C}(6) = 3.1213$. \emph{Status (2026):} AT2021 reports $3.102 \pm 0.032$, $\Delta = -0.61\%$, PASS.

\paragraph{F-T-6 (universal Boolean, $N\ge 2$).} For every $N\ge 2$ and the CR-canonical representative $K = \YKCR(N)$ per Table~\ref{tab:KASP}, no exceptional discrete eigenvalue exists in $(21/25, 1)$ on the privileged $\chi_*$-isotypic component of $L^2_{\mathrm{cusp}}(\YK)_{\chi_*}$. A single counterexample (an exceptional eigenvalue found at any anchor) falsifies the Transport Conjecture.

\smallskip

\noindent
The RMS deviation across F-T-1 to F-T-5 is $0.7\%$, comfortably inside the AT2021 continuum error bars $\sim 3\%$. No datapoint deviates by more than the $1\sigma$ AT2021 error bar from $\mathcal{C}(N)$. \emph{The universal Boolean F-T-6 is the cleanest single falsifier} and is the recommended primary test.

% =============================================================================
\section{Honest open-problem status (V2 dashboard)}\label{sec:status}
% =============================================================================

The Transport Conjecture (V2) is an \emph{open problem} with the following tier structure:
\begin{itemize}[leftmargin=2em]
   \item \textbf{At uniform $N$:} TIER 3 OPEN. Conjecture~\ref{conj:transport-formal-V2} is open; 4 candidate strategies (\S\ref{sec:strategies}) all open; the 2-Millennium-stack (M1) + (M2) scoping (\S\ref{sec:two-mill}) clarifies that any single proof strategy faces two Clay-level walls.
   \item \textbf{At $N=3$ specifically:} TIER 2 NUM operationally. The Boolean falsifier F-T-2 in \S\ref{sec:falsif} is computable in PARI/GP at $\leq \$50$ KVM budget, $\leq 3$ months elapsed time.
   \item \textbf{Empirically:} consistent with all known lattice data (AT2021, $0.7\%$ RMS over 6 datapoints). No counterexample known.
   \item \textbf{Strategically:} all four candidate proof strategies in \S\ref{sec:strategies} are themselves open programmes. The most tractable near-term option is (b) of \S\ref{ss:joint}: a Wightman positivity restriction to the surrogate observable.
\end{itemize}

We invite the community to attempt the conjecture via any of these (or new) strategies.

\subsection{Probability estimates (Bayesian, honest)}\label{ss:proba}

Per the cross-paper analysis \cite{Remondiere2026YMClosureV2}, the Bayesian decomposition is:
\begin{itemize}[leftmargin=2em]
  \item $P(\text{Transport TRUE} \mid \text{empirical RMS } 0.7\%) \approx 65$-$80\%$ (suggestive but not conclusive);
  \item $P(\text{PROVED within 5y} \mid \text{TRUE}) \approx 10$-$15\%$ (one candidate path: PARI numerical NUM-tier verification for SU(3));
  \item $P(\text{PROVED within 15y} \mid \text{TRUE}) \approx 25$-$40\%$;
  \item $P(\text{PROVED within 30y} \mid \text{TRUE}) \approx 50$-$70\%$.
\end{itemize}
Joint: $P(\text{Clay YM via Transport, 15y}) \approx 0.65 \times 0.30 + 0.80 \times 0.40 = 19.5\%$--$32\%$, median $\approx 26\%$.

\subsection{What this paper does \emph{not} claim}\label{ss:not-claim}

\begin{itemize}[leftmargin=2em]
   \item No proof. The conjecture is stated as an open problem.
   \item No claim of solution to the Yang--Mills Clay Millennium problem. The conjecture is concerned with a specific surrogate observable, and even at TIER 2 NUM for $N=3$ it does not establish constructive existence of 4D Yang--Mills.
   \item No claim that the RMS $0.7\%$ empirical match is a proof. Numerical agreement, however suggestive, is not equivalent to a structural argument.
\end{itemize}

% =============================================================================
\section{Acknowledgments}\label{sec:acks}
% =============================================================================

The author thanks the open mathematical literature for the foundational treatments used (Elstrodt--Grunewald--Mennicke, Sarnak, Bunke--Olbrich, Athenodorou--Teper, Chatterjee--Cao--Hairer--Shen, Cao--Nissim--Sheffield, Streater--Wightman). The work is the result of a structured arithmetic-gauge correspondence programme; full credit for the constituent inputs is given in the companion papers \cite{Remondiere2026Selberg, Remondiere2026CenterRank, Remondiere2026CR, Remondiere2026RouteB}.

The author thanks the DS-Bot-Vast.AI worker pool of the crossed-cosmos programme (Wave B 20 DS V4 Pro workers) for the convergent identification of the PARI Selberg-trace falsifier (10 workers) and for the Boolean $\lambda_1 = 1$ reformulation (worker 8, angle 9). The 2-Millennium-stack scoping (\S\ref{sec:two-mill}) is the convergent finding of 8 workers across angles 1, 3, 4, 13, 16, 17, 18, 20.

% =============================================================================
\bibliographystyle{alpha}
\bibliography{refs}
% =============================================================================

\end{document}
